\section{Allowing hypotheses on supplied lines}
\label{sec:boundaries}

Algorithm~\ref{alg:boundary} implements the reduction below.
In this section a supplied line returns one of three answers:
negative, zero, or positive. The \emph{zero}, or \emph{on-line},
answer means the hypothesis lies on the queried line, and its
answer class is $S_\ell^0=\{h\in S:f_\ell(h)=0\}$.
A \emph{ternary decision tree} has an outgoing edge for each
possible answer; a proper such tree retains only nonempty outcomes.
Leaf depths and the two cost objectives are defined exactly as
in the binary model.
Use the two virtual binary predicates
\[
 b_1(h)=\mathbf1[f_\ell(h)<0],\qquad
 b_2(h)=\mathbf1[f_\ell(h)\le0].
\]
Their pair of bits distinguishes the three outcomes.
Requesting either bit is implemented by querying the original line
once and storing its full sign. Both bits are then known. For analysis
we count each requested virtual bit as a separate binary query;
the actual number of supplied-line queries can only be smaller.

These binary labelings have planar off-boundary representations.
Any query whose answer is constant on the hypothesis set can be
discarded. For every remaining query, choose $\epsilon_\ell$ to be
half the minimum nonzero $|f_\ell(h)|$ on the finite hypothesis set.
This rational value is computable in polynomial time and has
polynomial encoding length. Then $b_1$ is represented
by $f_\ell(x)<-\epsilon_\ell$ and $b_2$ by
$f_\ell(x)<\epsilon_\ell$. Neither shifted boundary contains a
hypothesis. The shifts are used only for preprocessing;
their answers are obtained from the original supplied line.

\begin{symboltable}
$f_\ell(h)$ & Affine value whose sign is the answer for hypothesis $h$.\\
$b_1(h),\ b_2(h)$ & Virtual bits for $f_\ell(h)<0$ and $f_\ell(h)\le0$; together they encode the ternary sign.\\
$\epsilon_\ell$ & Positive shift smaller than every nonzero $|f_\ell(h)|$ for this finite state.
\end{symboltable}

\begin{algorithm}[H]
\caption{Transfer a binary algorithm to on-line outcomes}
\label{alg:boundary}
\KwIn{Finite hypotheses, supplied ternary lines, binary tree-construction routine $B$.}
\KwOut{An executable tree using the original supplied lines.}
\ForEach{nonconstant supplied line $\ell$}{
  $\epsilon_\ell\gets\frac12\min_{h:f_\ell(h)\ne0}|f_\ell(h)|$\;
  Represent virtual bits $b_1=\mathbf1[f_\ell(h)<0]$ and
  $b_2=\mathbf1[f_\ell(h)\le0]$ by the two shifted off-boundary lines\;
}
Construct a virtual binary tree with $B$\;
Execute a requested virtual bit by querying its original line,
caching the entire ternary sign; answer further requests from that cache\;
\KwRet the contracted original-line tree\;
\end{algorithm}

Lines constant on the whole hypothesis set are discarded before taking
the minimum nonzero affine value. Proposition~\ref{prop:boundary}
proves that the approximation ratio increases by at most two.


\begin{proposition}[Boundary reduction]\label{prop:boundary}
An $\alpha$-approximation for either binary objective gives a
$2\alpha$-approximation for the corresponding three-answer objective.
\end{proposition}
\begin{proof}
Simulate every query of an optimal original tree by its two virtual
bits. Each hypothesis uses at most twice as many binary queries,
so the virtual optimum is at most twice the original optimum
for either objective. Conversely, execute a constructed virtual
tree by original line queries and cached signs. Its actual depth
is no larger than its virtual depth on every hypothesis.
The two comparisons prove the factor-two transfer.
\end{proof}

\begin{corollary}[Constant approximations with on-line outcomes]
\label{cor:constant-boundaries}
With on-line answers, expected depth has a polynomial-time
approximation factor below $10.365952$ for arbitrary positive priors,
and deterministic maximum depth has a factor below $8.929931$.
\end{corollary}
\begin{proof}
Apply Proposition~\ref{prop:boundary} to
Theorems~\ref{thm:two-phase-expected} and~\ref{thm:poly-slot-worst}.
\end{proof}
